We ended the previous section by demonstrating that adjoining a root of an irreducible polynomial to a field $F$ produces an extension whose elements are polynomials in that root with a bounded degree. We focuse on studying the extensions generated by a single element. First, we define what it means to generate a field.

\begin{definition}[Generated Field] \leavevmode \\
    Let $K/F$ be an extension of fields and let $\alpha, \beta, \gamma, ... \in K$. Then the smallest subfield of $K$ containing both $F$ and $\alpha, \beta, \gamma, ...$ is denoted $F(\alpha, \beta, \gamma, ...)$ is called the field generated by $\alpha, \beta, \gamma, ...$ over $F$.
\end{definition}

\begin{definition}[Simple Extension, Primitive Element] \leavevmode \\
    If $K$ is generated by a single element $\alpha$ over $F$, $K = F(\alpha)$, then $K$ is said to be a simple extension of $F$ and $\alpha$ is called a primitive element of the extension.
\end{definition}

By adjoining a root $\alpha$ of an irreducible polynomial $p(x)$ to $F$, we get a simple extension $F(\alpha)$. Furthermore, this extension is isomorphic to the quotient $F[x]/(p(x))$.

\begin{theorem}
    Let $F$ be a field and let $p(x) \in F[x]$ be irreducible over $F$. Suppose $K$ is an extension field of $F$ containing a root $\alpha$ of $p(x)$. Then
    $$F(\alpha) \cong \frac{F[x]}{\left(p(x)\right)}$$
\end{theorem}

\begin{proof}
    Consider $\varphi: F[x] \to F[\alpha]$ where $a(x) \mapsto a(\alpha)$. Here, $F[\alpha]$ are polynomials in $\alpha$ with coefficients in $F$. Note that, by closure, $F[\alpha] \subseteq F(\alpha)$ and since $F \subseteq K$ and $\alpha \in K$, we have $F(\alpha) \subseteq K$. Since addition and multiplication are defined exactly the same in $F[x]$ and $F[\alpha]$, $\varphi$ is a ring homomorphism. Since $\alpha$ is a root of $p(x)$,
    \begin{align*}
        \varphi(p(x)) &= p(\alpha) = 0 \\
        \implies p(x) &\in \ker \varphi \\
        \implies \left(p(x)\right) &\subseteq \ker \varphi
    \end{align*}
    We use this to induce a well-defined homomorphism from $F[x]/(p(x))$ to $F(\alpha)$. Define $\mu: F[x]/(p(x)) \to F(\alpha)$ by $a(x) + (p(x)) \mapsto \varphi(\alpha)$. Note
    $$a(x) + (p(x)) = b(x) + (p(x)) \implies a(x) = b(x) + q(x)p(x)$$
    for some $q(x) \in F[x]$. So
    \begin{align*}
        \mu(a(x) + (p(x))) &= \mu(b(x) + q(x)\left(p(x) + (p(x))\right)) \\
        &= \varphi(b(x) + q(x)(x)) \\
        &= b(\alpha) + q(\alpha)p(\alpha) \\
        &= b(\alpha) \\
        &= \varphi(b(x)) \\
        &= \mu(b(x) + (p(x)))
    \end{align*}
    Since the domain of $\mu$ is a field,
    $$\ker \mu = \set{0} \text{ or } \ker \mu = \frac{F[x]}{(p(x))}$$
    However, $\mu(1+(p(x))) = 1 \neq 0$ so $\mu$ is not identically zero. It follows that $\ker \mu = \set{0}$ and $\mu$ is injective. The first isomorphism theorem then gives
    $$\frac{F[x]}{(p(x))}\cong \mu \left(\frac{F[x]}{(p(x))}\right) \subseteq F(\alpha)$$
    Noting that $\mu(x + (p(x))) = \alpha$ and $\mu(m + (p(x))) = m$ for all $m \in F$. Hence $\mu\left(F[x]/(p(x))\right)$ contains both $F$ and $\alpha$. Thus by definition $F(\alpha) \subseteq \mu\left(F[x]/(p(x))\right)$. Therefore, $\mu$ is onto and further it is an isomorphism.
    \qed
\end{proof}

\begin{corollary}
    Suppose in the previous theorem $\deg\left(p(x)\right) = n$. Then 
    $$F(\alpha) = \set{a_0 + a_1\alpha + a_2\alpha^2+...+a_{n-1}\alpha^{n-1} : a_0, ..., a_{n-1} \in F}$$
\end{corollary}

\begin{example}
    Consider $\alpha = \sqrt[3]{2}\in \R$. Then $\alpha$ is a root of $x^3-2$ which is irreducible over $\Q$ by Eisenstein's criterion. So
    \begin{align*}
        \Q(\sqrt[3]{2}) &= \set{a+b\sqrt[3]{2} + c\left(\sqrt[3]{2}\right)^2 : a, b, c, \in \Q} \\
        \frac{\Q[x]}{(x^3-2)} &\cong \Q(\sqrt[3]{2})
    \end{align*}
    Furthermore, $x^3-2$ has roots $\sqrt[3]{2}\omega$ and $\sqrt[3]{2}\omega^2$ where $\omega = -\frac{1}{2} + \frac{\sqrt{3}}{2}i.$
    $$\Q\left(\sqrt[3]{2}\omega\right) = \set{a+b\left(\sqrt[3]{2}\omega\right) + c\left(\sqrt[3]{2}\omega\right)^2 : a,b,c \in \Q} \cong \frac{\Q[x]}{(x^3 - 2)}$$
\end{example}